\documentclass[11pt,a4paper]{book}
\usepackage[amsbb]{mtpro2}
\usepackage[no-math,cm-default]{fontspec}
\usepackage{xunicode}
\usepackage{xgreek}
\defaultfontfeatures{Mapping=tex-text,Scale=MatchLowercase}
\setmainfont[Mapping=tex-text,Numbers=Lining,Scale=1.0,BoldFont={Nimbus Roman Bold}]{Nimbus Roman}
\usepackage[left=2.00cm, right=2.00cm, top=2.00cm, bottom=3.00cm]{geometry}
\usepackage{tikz}
\usepackage{tkz-euclide,tkz-fct}
\usepackage{wrapfig}
\usetkzobj{all}
\usepackage{calc}
\usepackage{systeme,regexpatch}
\usepackage[framemethod=TikZ]{mdframed}
\usepackage{adjustbox}
\makeatletter
\def\sysdelim#1#2{\def\SYS@delim@left{#1}\def\SYS@delim@right{#2}}
\sysdelim\{. % reinitialize
% patch the internal command to use
% \LEFTRIGHT<left delim><right delim>{<system>}
% instead of \left<left delim<system>\right<right delim>
\regexpatchcmd\SYS@systeme@iii
  {\cB.\c{SYS@delim@left}(.*)\c{SYS@delim@right}\cE.}
  {\c{SYS@MT@LEFTRIGHT}\cB\{\1\cE\}}
  {}{}
\def\SYS@MT@LEFTRIGHT{%
  \expandafter\expandafter\expandafter\LEFTRIGHT
  \expandafter\SYS@delim@left\SYS@delim@right}
\makeatother
\newcommand{\synt}[2]{{\scriptsize \begin{matrix}
\times#1\\\\ \times#2
\end{matrix}}}
\usepackage{graphicx}
\usepackage{multicol}
\usepackage{multirow}
\usepackage{enumitem}
\usepackage{tabularx}
\newcommand{\xr}[1]{\textcolor{red!80!black}{#1}}
\setlist[enumerate]{itemsep=0mm}
\begin{document}
\chapter{Εργασίες - Μελλοντικά project}
\begin{enumerate}[label=\arabic*,itemsep=0mm]
\begin{multicols}{2}
\item Βιβλία Μαθηματικών
\item Βάση Δεδομένων - Ασκήσεων
\item Προγράμματα για βάση δεδομένων - Texomaker
\item Σημειώσεις
\item Τυπολόγιο
\item Χρονοδιάγραμμα ύλης
\item Πίνακες ύλης
\item Φυλλάδια ασκήσεων
\item Διαγωνίσματα
\item Πλάνο διδασκαλίας
\item Διαφάνειες
\item Σημαντικές ασκήσεις για εξετάσεις
\item Σελίδα Ιντερνετ - mathworld.gr
\item \LaTeX\ online
\item Blog
\item Διαδραστικά τεστ με βραβεία,\\επίπεδα, βαθμολογία
\item Ιστορία - Βιογραφίες - Γεγονότα
\item Παλιά βιβλία - περιοδικά
\item Λυμένα θέματα εξετάσεων
\item Εργασίες σε διάφορα θέματα
\item Προβλήματα από όλες τις επιστήμες
\item Μαθηματικοί διαγωνισμοί - LaTeX
\item Mathematica
\item Matlab - Mupad
\item \LaTeX\ Templates - Classes - Packages
\item Διάφοροι πίνακες
\item Γρίφοι παιχνίδια
\item Δακτυλογραφήσεις κειμένων - εργασιών σε \LaTeX
\item Βίντεο Μαθήματα
\item Συμβουλές προγραμματισμού
\item Διαφάνειες Javascript
\item Αφίσες / Είδη δώρων / Διαφημιστικά / Ημερολόγια / Μαθηματικά ρολόγια
\item Τηλεφωνικός κατάλογος μαθητών
\item Φυλλάδιο ύλης με βαθμούς
\item Διαφήμιση
\item Βάση μαθητών - Access
\item Αρχεία Maple
\item Αρχεία Geogebra
\end{multicols}
\end{enumerate}

\newpage
\chapter{Βιβλία}
\section{Περιεχόμενα}
\begin{multicols}{2}
\begin{enumerate}[itemsep=0mm]
\item Μαθηματικά
\item Βοηθήματα Γυμνασίου
\item Βοηθήματα Λυκείου
\item Τα εκτός ύλης
\item Διαφορικές Εξισώσεις
\item Σχολικό τυπολόγιο
\item Εφαρμοσμένα Μαθ. Γυμνασίου - Λυκείου
\item Βιβλία πανεπιστημίου
\item Προετοιμασία Γ΄ Λυκείου
\item Μαθηματική εγκυκλοπαίδεια - Λεξικό
\item Λεξικό 
\end{enumerate}
\end{multicols}
\section{Βοηθήματα Γυμνασίου - Λυκείου}
\subsection{Περιεχόμενα}
\begin{enumerate}
\item Θεωρία με ορισμούς, θεωρήματα, παρατηρήσεις, συμβουλές. Η θεωρία με λίγη περιγραφή (όχι μόνο λίστα).
\item Μέθοδοι με λυμένες ασκήσεις. Άλυτες ασκήσεις. Μετά από τις μεθόδους, γενικές ασκήσεις με διάγραμμα λύσης. Παραδείγματα μετά από κάθε μέθοδο. Άλυτες ασκήσεις μοιρασμένες ανά μέθοδο και γενικές. Ασκήσεις με πραγματικά στοιχεία.
\item Διαγωνίσματα. Ύστερα από κάθε κεφάλαιο αλλά και στο τέλος επαναληπτικά. 
\item Επαναληπτικές ασκήσεις. Συνδυασμός κομματιών της ύλης και προηγούμενων τάξεων. Στο τέλος. Προβλήματα.
\item Χρήσιμοι πίνακες. Παράρτημα.
\item Βιβλιογραφία.
\item Ιστορικά στοιχεία. Μαθηματικοί.
\item Θέματα εξετάσεων μέσα στις ασκήσεις (Γ΄ Λυκείου). Λυμένα θέματα πανελλαδικών.
\item Απαντήσεις ασκήσεων στο τέλος ή ξεχωριστό βιβλίο.
\item Τυπολόγιο στο τέλος κάθε κεφαλαίου.
\item Παραπομπές στη σελίδα.
\end{enumerate}
\subsection{Δομή}
\begin{enumerate}
\item Εσώφυλλο με διάγραμμα, μετά αφιέρωση, πρόλογος.
\item Πίνακας με συμβολισμούς.
\item Περιεχόμενα
\item Κεφάλαιο
\begin{enumerate}
\item Παράγραφος
\begin{enumerate}
\item Ορισμοί
\item Θεωρήματα, ιδιότητες, πορίσματα κτλ.
\item Μέθοδοι με λυμένα παραδείγματα αμέσως μετά.
\item Συμπληρωματικές ασκήσεις.
\item Άλυτες ασκήσεις.
\end{enumerate}
\item Διαγωνίσματα
\item Τυπολόγιο
\end{enumerate}
\item Παράρτημα
\item Πίνακες
\item Απαντήσεις ασκήσεων
\item Βιβλιογραφία
\end{enumerate}
\section{Τα εκτός ύλης}
\subsection{Περιεχόμενα}
\begin{enumerate}
\item Θεωρία με πλήρη περιγραφή, απόδειξη θεωρημάτων, πορίσματα, παρατηρήσεις.
\item Μέθοδοι με παράδειγμα αμέσως μετά.
\item Άλυτες ασκήσεις.
\item Ιστορικά στοιχεία.
\item Παραπομπές στη σελίδα.
\item Απαντήσεις ασκήσεων.
\item Βιβλιογραφία.
\item Παράρτημα;
\end{enumerate}
\subsection{Δομή}
\begin{enumerate}
\item Εσώφυλλο με διάγραμμα, μετά αφιέρωση, πρόλογος.
\item Πίνακας με συμβολισμούς.
\item Περιεχόμενα
\item Κεφάλαιο
\begin{enumerate}
\item Παράγραφος
\begin{enumerate}
\item Ορισμοί
\item Θεωρήματα, ιδιότητες, πορίσματα κτλ.
\item Μέθοδοι με λυμένα παραδείγματα αμέσως μετά.
\item Συμπληρωματικές ασκήσεις.
\item Άλυτες ασκήσεις.
\end{enumerate}
\end{enumerate}
\item Παράρτημα
\item Πίνακες
\item Απαντήσεις ασκήσεων
\item Βιβλιογραφία
\end{enumerate}
\section{Διαφορικές Εξισώσεις}
\subsection{Περιεχόμενα}
\begin{enumerate}
\item Θεωρία του βιβλίου και γενική. Λίστα.
\item Μεθοδολογία.
\item Λύσεις ασκήσεων βιβλίου.
\item Λύσεις ασκήσεων φυλλαδίου.
\item Πίνακες.
\item Βιβλιογραφία.
\end{enumerate}
\section{Σχολικό τυπολόγιο}
\section{Εφαρμοσμένα μαθηματικά Γυμνασίου - Λυκείου}
\subsection{Κλάδοι}
\begin{multicols}{2}
\begin{enumerate}
\item Φυσική
\item Χημεία
\item Βιολογία
\item Ιατρική
\item Πληροφορική
\item Παιχνίδια
\item Φωτογραφία - Οπτική
\item Αθλητισμός
\item Αυτοκίνητα - Μηχανές
\item Μουσική - μουσικά όργανα
\item Σπίτι - Κτίρια - κατασκευές - μνημεία
\item Τοποθεσίες
\item Αστροφυσική
\item Εμπόριο
\item Οικονομία
\item Πλοία
\item Φυσικά φαινόμενα (σεισμοί, καιρός,...)
\item Αεροπλάνα
\item Σχεδίαση με καμπύλες Beizer (Η/Υ)
\item Τηλεπικοινωνίες
\item Αντικείμενα - Συσκευές
\item Γεωγραφία - Χάρτες
\item 
\end{enumerate}
\end{multicols}
\chapter{Βάση Δεδομένων}
\section{Περιεχόμενα}
\begin{multicols}{2}
\begin{enumerate}[itemsep=0mm]
\item Ορισμοί
\item Θεωρήματα
\item Άλυτες ασκήσεις
\item Λύσεις ασκήσεων
\item Μέθοδοι για λύση ασκήσεων
\item Λυμένα παραδείγματα
\item Θέματα διαγωνισμάτων ανά ενότητα
\item Λύσεις θεμάτων ανά ενότητα
\item Συνδυαστικά θέματα
\item Λύσεις συνδυαστικών θεμάτων
\item Πίνακες
\item Σχήματα
\end{enumerate}
\end{multicols}
\section{Κωδικοποίηση}
\subsection{Κωδικοποίηση φακέλων}\label{par:KwdFakelwn}
\begin{itemize}[leftmargin=4mm,itemsep=0mm]
\item Για το φάκελο "Μάθημα"
\begin{multicols}{2}
\begin{enumerate}
\item Άλγεβρα $ \to $ \xr{Algebra}
\item Γεωμετρία $ \to $ \xr{Gewmetria}
\item Ανάλυση $ \to $ \xr{Analysh}
\item Πιθανότητες - Στατιστική $ \to $ \xr{PithanStat}
\end{enumerate}
\end{multicols}
\item Για τους φακέλους "Ενότητα" τα ονόματα είναι καταχωρημένα στο αρχείο "\xr{Kwdikoi.csv}" και θα αντλούνται από το πρόγραμμα. Οι κωδικοί θα προκύπτουν από σύντομη ονομασία των ενοτήτων.
\item Για τους φακέλους "Είδος Άσκησης" τα ονόματα είναι καταχωρημένα στο αρχείο "\xr{Kwdikoi.csv}" και θα αντλούνται από το πρόγραμμα. Οι κωδικοί θα προκύπτουν από σύντομη ονομασία του είδους άσκησης.
\end{itemize}
\subsection{Κωδικοποίηση ορισμών}
Στην κωδικοποίηση των ορισμών, ο κωδικός του ορισμού θα προκύπτει από την πλήρη ονομασία του. Π.χ. για τον ορισμό της τετραγωνικής ρίζας ο κωδικός θα είναι \xr{TetragwnikhRiza}. Αν χρειαστεί ο ίδιος ορισμός να δοθεί δύο φορές τότε προσθέτουμε αριθμό στο τέλος. Π.χ. \xr{TetragwnikhRiza1} για το γυμνάσιο και \xr{TetragwnikhRiza2} για το λύκειο. Μετά το "Μάθημα" θα υπάρχει ο κωδικός "\xr{Or}" για τη δήλωση του ορισμού.
\begin{center}
<Όνομα>-<Μάθημα>-Or-<Κωδικό Όνομα Ορισμού>.tex\\
\xr{S}-Algebra-\xr{Or}-TetragwnikhRiza.tex
\end{center}
\subsection{Κωδικοποίηση θεωρημάτων}
Στην κωδικοποίηση των θεωρημάτων, ο κωδικός του θεωρήματος θα προκύπτει από την πλήρη ονομασία του. Π.χ. για τις ιδιότητες των ριζών ο κωδικός θα είναι \xr{IdiothtesRizwn}. Αν χρειαστεί ο ίδιος ορισμός να δοθεί δύο φορές τότε προσθέτουμε αριθμό στο τέλος. Μετά το "Μάθημα" θα υπάρχει ο κωδικός "\xr{Th}" για τη δήλωση του θεωρήματος.
\begin{center}
<Όνομα>-<Μάθημα>-Th-<Κωδικό Όνομα Θεωρήματος>.tex\\
\xr{S}-Algebra-\xr{Th}-IdiothtesRizwn.tex
\end{center}
\subsection{Κωδικοποίηση Άλυτων Ασκήσεων}
Για την κωδικοποίηση των άλυτων ασκήσεων χρησιμοποιούμε τον κωδικό ενότητας και τον κωδικό για το είδος άσκησης όπως αυτά ορίστηκαν στην παράγραφο \ref{par:KwdFakelwn}. Ο κωδικός \xr{AA} αντιστοιχεί στο "Άλυτη άσκηση" και ακολουθείται από ένα διψήφιο αριθμό.
\begin{center}
<Όνομα>-<Ενότητα>-<Είδος Άσκησης>-AA/Αριθμός.tex\\
\xr{S}-GrSys-\xr{MetAn}-AA01.tex
\end{center}
\subsection{Κωδικοποίηση Λύσεων Ασκήσεων}
Για την κωδικοποίηση των Λύσεων των ασκήσεων χρησιμοποιούμε το ίδιο μοτίβο με την κωδικοποίηση των άλυτων με τη διαφορά ότι χρησιμοποιούμε τον κωδικό \xr{LA} που αντιστοιχεί στο "Λύση άσκησης" και ακολουθείται από ένα διψήφιο αριθμό.
\begin{center}
<Όνομα>-<Ενότητα>-<Είδος Άσκησης>-LA/Αριθμός.tex\\
\xr{S}-GrSys-\xr{MetAn}-LA01.tex
\end{center}
\subsection{Κωδικοποίηση Μεθόδων}
Στην κωδικοποίηση των μεθόδων, ο κωδικός της μεθόδου θα προκύπτει από την πλήρη ονομασία της. Π.χ. για τη μέθοδο της αντικατάστασης ο κωδικός θα είναι \xr{MethodosAntikatastashs}. Μετά την "Ενότητα" θα υπάρχει ο κωδικός "\xr{Meth}" για τη δήλωση της μεθόδου.
\begin{center}
<Όνομα>-<Ενότητα>-Meth-<Όνομα Μεθόδου>.tex\\
\xr{S}-GrSys-\xr{Meth}-MethodosAntikatastashs.tex
\end{center}
\subsection{Κωδικοποίηση Λυμένων Παραδειγμάτων}
Για την κωδικοποίηση των λυμένων παραδειγμάτων χρησιμοποιούμε το ίδιο μοτίβο με την κωδικοποίηση των μεθόδων, με τη διαφορά ότι χρησιμοποιούμε τον κωδικό \xr{LP} που αντιστοιχεί στο "Λυμένο παράδειγμα", ακολουθούμενο από ένα διψήφιο αριθμό.
\begin{center}
<Όνομα>-<Ενότητα>-LP-<Όνομα Μεθόδου>/Αριθμός.tex\\
\xr{S}-GrSys-\xr{LP}-MethodosAntikatastashs01.tex
\end{center}
\subsection{Κωδικοποίηση Θεμάτων Ενότητας}
Στην κωδικοποίηση των άλυτων θεμάτων ενότητας, ο κωδικός ενότητας είναι όπως ορίστηκε προηγουμένως. Τα είδη των θεμάτων καθώς και οι κωδικές ονομασίες τους είναι:

\begin{center}
<Όνομα>-<Ενότητα>-<Είδος Θέματος>-AT/Αριθμός.tex\\
\xr{S}-GrSys-\xr{Thewria}-AT01.tex
\end{center}
\subsection{Κωδικοποίηση Λύσεων Θεμάτων Ενότητας}
\begin{center}
<Όνομα>-<Ενότητα>-<Είδος Θέματος>-LT/Αριθμός.tex\\
\xr{S}-GrSys-\xr{Thewria}-LT01.tex
\end{center}
\subsection{Κωδικοποίηση Συνδυαστικών Θεμάτων}
\begin{center}
<Όνομα>-<Ένωση Κωδικών>-SyndLT/Αριθμός.tex\\
\xr{S}-MonotKyrtEfapt-\xr{SyndAT01}.tex
\end{center}
\subsection{Κωδικοποίηση Λύσεων Συνδυαστικών Θεμάτων}
\begin{center}
<Όνομα>-<Ένωση Κωδικών>-SyndLT/Αριθμός.tex\\
\xr{S}-MonotKyrtEfapt-\xr{SyndLT01}.tex
\end{center}
\subsection{Κωδικοποίηση Πινάκων}
Για την κωδικοποίηση των πινάκων, ο κωδικός του πίνακα προκύπτει από την πλήρη ονομασία του. Για παράδειγμα ο πίνακας με τις ιδιότητες των ριζών θα έχει κωδικό "\xr{IdiothtesRizwn}". Μετά το "Μάθημα" χρησιμοποιούμε των κωδικό "\xr{Pin}" που αντιστοιχεί στη λέξη Πίνακας.
\begin{center}
<Όνομα>-<Μάθημα>-Pin<Κωδικό Όνομα Πίνακα>.tex\\
\xr{S}-Algebra-\xr{Pin}-IdiothtesRizwn.tex
\end{center}
\subsection{Κωδικοποίηση Σχημάτων}
Για την κωδικοποίηση των πινάκων, ο κωδικός του πίνακα προκύπτει από την πλήρη ονομασία του. Για παράδειγμα ο τριγωνομετρικός κύκλος θα έχει κωδικό "\xr{TrigonometrikosKyklos}". Μετά το "Μάθημα" χρησιμοποιούμε των κωδικό "\xr{Sx}" που αντιστοιχεί στη λέξη Σχήμα.
\begin{center}
<Όνομα>-<Μάθημα>-Sx-<Κωδικό Όνομα Σχήματος>.tex\\
\xr{S}-Algebra-\xr{Sx}-TrigonometrikosKyklos.tex
\end{center}
\chapter{Πρόγραμμα για βάση δεδομένων - Texomaker}
\section{Περιγραφή - Περιεχόμενα}
\subsection{Περιεχόμενα}

\subsection{Περιγραφή}
\chapter{Διαγωνίσματα}
\section{Τύποι διαγωνισμάτων}
\begin{multicols}{2}
\begin{enumerate}[itemsep=0mm]
\item Τύπου Α - Θέματα θεωρίας
\item Τύπου Β - Κλασικό διαγώνισμα
\item Τύπου Γ - Συνδυαστικά θέματα
\item Τύπου Δ - Προβλήματα
\item Τύπου Ε - Διαγνωστικά τεστ
\item Τύπου ΣΤ - Θέματα σχολικού βιβλίου
\item Τύπου Ζ - Θέματα εξετάσεων
\item Τύπου Η - Θέματα αυξημένης δυσκολίας
\item Τύπου Θ - Μαθηματικοί διαγωνισμοί
\end{enumerate}
\end{multicols}
\section{Δομή Διαγωνισμάτων - Περιεχόμενα θεμάτων}
Τα διαγωνίσματα θα έχουν μια από τις ακόλουθες δομές:
\begin{enumerate}
\item \textit{Τυπική Δομή}\\
Για το Γυμνάσιο και το Λύκειο το διαγώνισμα θα περιέχει $ 4 $ θέματα διαβαθμισμένης δυσκολίας. Μερικά από τα θέματα θα αφορούν μια ενότητα της εξεταστέας ύλης και τα υπόλοιπα θα είναι συνδυαστικά, εκτός και αν ορίζεται αλλιώς για κάποιο τύπο διαγωνίσματος.
\item \textit{Σχολική Δομή}
\begin{itemize}[leftmargin=0mm]
\item Για το Λύκειο το διαγώνισμα περιέχει $4$ θέματα. Μερικά από τα θέματα θα αφορούν μια ενότητα της εξεταστέας ύλης και τα υπόλοιπα θα είναι συνδυαστικά, εκτός και αν ορίζεται αλλιώς για κάποιο τύπο διαγωνίσματος.
\item Για το Γυμνάσιο το διαγώνισμα περιέχει $ 5 $ θέματα μοιρασμένα σε $ 2 $ ομάδες, $ 2 $ θέματα η ομάδα $ A $ και $ 3 $ θέματα η ομάδα $ B $.Ο μαθητής επιλέγει $ 3 $ θέματα, ένα από την ομάδα $ A $ και δύο από τη $ B $. Τα θέματα θα είναι κάποια από μια ενότητα της ύλης και τα υπόλοιπα συνδυαστικά, εκτός και αν ορίζεται αλλιώς.
\end{itemize}
\end{enumerate}
\section{Τύπου Α Θεωρία}
\noindent
Τα διαγωνίσματα "Τύπου A" έχουν είτε \textit{Τυπική δομή} είτε \textit{Σχολική δομή} και αποτελούνται αποκλειστικά από θέματα πάνω στη θεωρία. Τα είδη ασκήσεων που περιέχονται είναι
\begin{multicols}{2}
\begin{itemize}[itemsep=0mm]
\item Διατύπωση Ορισμού - Πρότασης
\item Απόδειξη θεωρήματος
\item Ερωτήσεις Σωστό - Λάθος
\item Αντιστοίχηση
\item Πολλαπλής επιλογής
\item Αντιπαραδείγματα
\item Διόρθωση λάθους
\item Επιλογή θεωρίας για επίλυση άσκησης
\item Συμπλήρωση κενού
\end{itemize}
\end{multicols}
\section{Τύπου Β - Κλασικό Διαγώνισμα - Οδηγίες Υπουργείου}
\noindent
Τα διαγωνίσματα "Τύπου Β" περιέχουν θέματα σύμφωνα με τις οδηγίες του υπουργείου. Συγκεκριμένα έχουν \textit{Σχολική δομή} όπου
\begin{itemize}
\item για το γυμνάσιο περιέχουν 2 θέματα θεωρίας με κάθε θέμα να καλύπτει ένα αντικείμενο της εξεταστέας ύλης καθώς και 3 ασκήσεις ίδιου τύπου.
\item για το λύκειο περιέχει 4 θέματα. Συγκεκριμένα
\begin{enumerate}
\item το πρώτο θέμα αποτελείται από δύο μέρη. Το πρώτο μέρος περιέχει πέντε (05) ερωτήσεις αντικειμενικού τύπου (πολλαπλής επιλογής, Σωστού - Λάθους, αντιστοίχισης). Στο δεύτερο μέρος ζητείται η απόδειξη μίας απλής πρότασης (ιδιότητας, λήμματος, θεωρήματος ή πορίσματος), που είναι αποδεδειγμένη στο σχολικό εγχειρίδιο.
\item Το δεύτερο θέμα αποτελείται από μία άσκηση που είναι εφαρμογή ορισμών, αλγορίθμων ή προτάσεων (ιδιοτήτων, θεωρημάτων, πορισμάτων).
\item Το τρίτο θέμα αποτελείται από μία άσκηση που απαιτεί από τον μαθητή ικανότητα συνδυασμού και σύνθεσης εννοιών και αποδεικτικών ή υπολογιστικών διαδικασιών.
\item Το τέταρτο θέμα αποτελείται από μία άσκηση ή ένα πρόβλημα που η λύση του απαιτεί από τον μαθητή ικανότητες συνδυασμού και σύνθεσης γνώσεων, αλλά και την ανάληψη πρωτοβουλιών για την ανάπτυξη στρατηγικών επίλυσής του.
\end{enumerate}
\end{itemize}
\section{Τύπου Γ - Συνδυαστικά Θέματα}
\noindent
Τα διαγωνίσματα "Τύπου Γ" έχουν \textit{Τυπική δομή} και αποτελούνται από $4$ συνδυαστικά θέματα πάνω στην εξεταστέα ύλη της χρονιάς. Κάθε θέμα περιλαμβάνει διάφορα τμήματα της ύλης και αποτελείται από διαδοχικά ερωτήματα. Βαθμολογείται με 4/20 μονάδες για το γυμνάσιο και 25/100 μονάδες για το λύκειο.
\section{Τύπου Δ - Προβλήματα}
Τα διαγωνίσματα "Τύπου Δ" έχουν είτε \textit{Τυπική δομή} είτε \textit{Σχολική δομή} και αποτελούνται αποκλειστικά από προβλήματα.
\section{Τύπου Ε - Διαγνωστικά Τεστ}
\noindent
Τα διαγωνίσματα "Τύπου Ε" είναι διαγνωστικά τεστ που διεξάγονται στην αρχή της σχολικής χρονιάς με σκοπό την αξιολόγηση των γνώσεων του μαθητή στην περσινή ύλη. Καλύπτουν όσο το δυνατόν περισσότερο μέρος της εξεταστέας ύλης κατά κύριο λόγο της περασμένης χρονιάς αλλά και των προηγούμενων τάξεων. 
\section{Τύπου ΣΤ - Θέματα Σχολικού Βιβλίου}
\noindent
Τα διαγωνίσματα "Τύπου ΣΤ" έχουν \textit{Σχολική δομή} και περιέχουν ασκήσεις επιλεγμένες από το σχολικό βιβλίο.
\section{Τύπου Ζ - Θέματα Εξετάσεων}
\noindent
Τα διαγωνίσματα "Τύπου Ζ" είναι τα διαγωνίσματα στα οποία έχουν εξεταστεί οι μαθητές στις σχολικές εξετάσεις.
\begin{itemize}[itemsep=0mm]
\item Για τις τάξεις του Γυμνασίου καθώς και την Α΄ και Β΄ Λυκείου τα διαγωνίσματα θα είναι επιλεγμένα από διάφορα σχολεία της Ελλάδας.
\item Για τη Γ΄ Λυκείου τα διαγωνίσματα θα είναι θέματα πανελλαδικών των περασμένων ετών.
\end{itemize}
\section{Τύπου Η - Θέματα Αυξημένης Δυσκολίας}
Τα διαγωνίσματα "Τύπου Η" έχουν είτε \textit{Τυπική δομή} είτε \textit{Σχολική δομή} και αποτελούνται από θέματα όλων των ειδών:
\begin{itemize}[itemsep=0mm]
\item Θέματα θεωρίας
\item Υπολογιστικές ασκήσεις
\item Συνδυαστικά θέματα
\item Προβλήματα
\end{itemize}
Τα θέματα είναι αυξημένης δυσκολίας και απευθύνονται σε μαθητές που έχουν καλύψει αρκετά καλά τις βασικές ενότητες της εξεταστέας ύλης και θέλουν να εξεταστούν πάνω σε πιο απαιτητικά θέματα.
\section{Τύπου Θ - Μαθηματικοί διαγωνισμοί}
\noindent
Τα διαγωνίσματα "Τύπου Θ" είναι θέματα Μαθηματικών διαγωνισμών που διοργανώνονται από την Ελληνική Μαθηματική Εταιρία. Αφορούν τους μαθηματικούς διαγωνισμούς "Ευκλείδης", "Θαλής", "Αρχιμήδης "..... όλων των τάξεων Γυμνασίου και Λυκείου καθώς και θέματα μαθηματικών διαγωνισμών από παραρτήματα της Ε.Μ.Ε. στις διάφορες πόλεις της Ελλάδας.
\end{document}

